\begin{idstatement}{C.3.9}{Proved exclusion of secondary Kodaira general fibres}
\label{st:C.3.9}
A compact complex threefold with $b_1(Y)=b_2(Y)=0$ and
$a(Y)=1$ cannot have a holomorphic algebraic reduction
$f:Y\to\PP^1$ whose general fibre is a secondary Kodaira
surface. Rational Betti-number vanishing suffices; no integral
torsion-freeness of the fibre is assumed.
\end{idstatement}
\begin{proof}
Suppose that the smooth general fibres $S_c$ are secondary
Kodaira surfaces. They have $h^1(\OO_{S_c})=b_1(S_c)=1$.
The nonzero integral transgression section of D.1.9 has
coefficient image $\beta\in H^1(Y,\OO_Y)$ and gives a
nowhere-zero section of $\mathscr E=R^1f_*\OO_Y$.
This coherent sheaf is a line bundle by E4. Thus
$\mathscr E=\OO_{\PP^1}$ and
$H^1(Y,\OO_Y)=\C\beta$. For every $z\in\C$ set
\[
 M_z=\exp(2\pi\mathrm i z\beta)\in\Pic^0(Y).
\tag{C.SK.1}
\]
Its restriction satisfies $M_{z+1}|_{S_c}\simeq M_z|_{S_c}$,
because $\beta|_{S_c}$ is the coefficient image of an
integral class. It is nonzero as a tangent vector of
$\Pic^0(S_c)$ by D.1.7. The exponential sequence therefore
identifies $\Pic^0(S_c)$ with $\C^*$ and makes (C.SK.1), on
each smooth fibre, a nonconstant homomorphism from
$\C/\Z$ to this group.

We specify the primary cover and the invariant used below.
Let $m>1$ be the order of $K_{S_c}$, and let
\[
 p_c:P_c\longrightarrow S_c
\]
be its connected cyclic canonical cover, with deck group
$G=\mu_m$. It is finite \emph{\'etale}. Its canonical bundle
is trivial, and it is non-K\"ahler, because a finite
\emph{\'etale} quotient of a K\"ahler surface is K\"ahler
by averaging a K\"ahler form. The classification of
surfaces with trivial canonical bundle consequently makes
$P_c$ a primary Kodaira surface. Write
$\pi_c:P_c\to E_c$ for its Albanese elliptic bundle, and
$d>0$ for its topological degree. These covers exist in
holomorphic families locally on the smooth base: the
evaluation of a local frame of $f_*K_{Y/\PP^1}^{m}$
trivializes $K_{Y/\PP^1}^{m}$, since it is nowhere zero on
each fibre, and taking its cyclic root gives the required
finite \emph{\'etale} covering family. The integers $m,d$
are locally constant and hence constant on the connected
smooth base.

The exact primary-surface invariant we use is
\[
 Q_c=\Pic^0(P_c)/\pi_c^*\Pic^0(E_c)\simeq\C^*.
\tag{C.SK.2}
\]
The Picard pullback in the denominator is injective:
$\pi_{c*}\OO_{P_c}=\OO_{E_c}$ and the projection formula
show that a line with trivial pullback is trivial. For
any line $L_c$ of degree $d$ on $E_c$, the pullback
$\pi_c^*L_c$ belongs to $\Pic^0(P_c)$, and its class in
$Q_c$ is independent of $L_c$. Choose the group
isomorphism $q_c:Q_c\to\C^*$ that places this
positive-degree class in the punctured unit disc, and
denote its value by $\alpha(c)$.
This fixes the possible inversion ambiguity of the
group coordinate. The descriptions in
\cite[Theorem~4 and Remark~1]{Borcea1984} and
\cite[\S1, pp.~2--3]{SS2002} identify $\alpha$ as the
intrinsic holomorphic punctured-disc coordinate of the
moduli of primary Kodaira surfaces. More explicitly,
$P_c$ is the associated elliptic bundle of a line of
degree $d$ with fibre $\C^*/\langle\alpha(c)\rangle$;
quotienting by $\pi_c^*\Pic^0(E_c)$ removes the choice
of that line within its degree. In a holomorphic
family the local integral marking and the coefficient
map give this quotient a holomorphic $\C^*$ coordinate,
and the chosen positive-degree orientation is locally
constant. This is also immediate in Borcea's complete
period family. Thus $\alpha(c)$ is holomorphic.

The locally constructed canonical covers have the same
isomorphism class on overlaps, so this intrinsic
function glues on the smooth base. It has
$0<|\alpha(c)|<1$. It extends across the finitely many
omitted values by the removable singularity theorem
and is constant on $\PP^1$. Write its nonzero constant
value as $\alpha$.

Consider the exact homomorphism of connected Picard
groups
\[
 \Pic^0(S_c)\xrightarrow{\ p_c^*\ }\Pic^0(P_c)
          \longrightarrow Q_c\xrightarrow{\ q_c\ }\C^*.
\tag{C.SK.3}
\]
Its differential is nonzero. Indeed finite-cover trace
identifies $H^1(S_c,\OO_{S_c})$ with the $G$-invariant
line in $H^1(P_c,\OO_{P_c})$. The other line,
$\pi_c^*H^1(E_c,\OO_{E_c})$, has nontrivial $G$ action:
otherwise the unique holomorphic one-form of $P_c$
would descend to $S_c$, whereas a secondary Kodaira
surface has no holomorphic one-form. Hence the invariant
line has zero intersection with the denominator's
tangent line in (C.SK.2). This proves the claim.
Consequently (C.SK.3) is a nonzero isogeny of $\C^*$
groups. Composing it with (C.SK.1) gives
\[
 q_c\bigl([p_c^*(M_z|_{S_c})]\bigr)
       =\exp(2\pi\mathrm i k z)
       \qquad(k\in\Z\setminus\{0\}).
\tag{C.SK.4}
\]
Here integrality follows from the period $z\mapsto z+1$;
every holomorphic group homomorphism $\C^*\to\C^*$ is
an integral power. The integer $k$ is locally constant
in the holomorphic families just described, and is
therefore a single integer on the connected smooth base.

Let $\rho_c:S_c\to\PP^1$ be its elliptic algebraic
reduction. The induced finite map $r_c:E_c\to\PP^1$
satisfies $\rho_c p_c=r_c\pi_c$. Write $h=\deg r_c$;
this positive integer is constant. Choose one positive
integer $N$ killing the torsion group $H^2(S_c;\Z)$
and satisfying $d\mid hN$. This choice is uniform
because the smooth fibres are differentiably locally
trivial. Then
$A_c=\rho_c^*\OO_{\PP^1}(N)$ belongs to $\Pic^0(S_c)$,
and (C.SK.2) gives
\[
 q_c([p_c^*A_c])=\alpha^{hN/d}.
\tag{C.SK.5}
\]
This identity is fibrewise and uses only degrees; it
does not require the divisors $\rho_c^{-1}(x)$ to
extend across a singular fibre of $f$.

Choose $z_0\in\C$ with
$\exp(2\pi\mathrm i k z_0)=\alpha^{hN/d}$ and put
$M=M_{z_0}$. By (C.SK.4)--(C.SK.5), the difference of
$p_c^*(M|_{S_c})$ and $p_c^*A_c$ is
$\pi_c^*R_c$ for some $R_c\in\Pic^0(E_c)$.
Both of the first two lines come from $S_c$ and are
$G$-invariant. Injectivity of $\pi_c^*$ implies that
$R_c$ is fixed by the induced nontrivial linear action
on $\Pic^0(E_c)$. An elliptic-curve linear automorphism
of finite order other than one has order
$r\in\{2,3,4,6\}$. If it fixes a point $R_c$, then
\[
 rR_c=(1+\zeta+\cdots+\zeta^{r-1})R_c=0
\]
in the elliptic group. Thus $R_c^{\otimes12}$ is
trivial. The norm of the finite cover satisfies
$\operatorname{Nm}_{p_c}(p_c^*L)=L^{\otimes m}$.
Applying it to the difference now proves the exact
line-bundle identity
\[
 M^{\otimes12m}|_{S_c}
       \simeq \rho_c^*\OO_{\PP^1}(12mN)
\tag{C.SK.6}
\]
for every smooth general fibre.

The right side has at least two linearly independent
sections, by pullback from $\PP^1$. Coherent generic
base change gives
$\operatorname{rank}f_*M^{\otimes12m}\ge2$.
Its relative linear system has a positive-dimensional
image on a general fibre and therefore a surface
image in the projective bundle over $\PP^1$.
Pullback of meromorphic functions from this image
gives $a(Y)\ge2$, the desired contradiction.
This proves the exclusion without the arbitrary-twist
cohomological-flatness assertion used in the printed
Kodaira exclusion.
\end{proof}


% C.3.10 from part_C1_C3.tex

